\documentclass[10pt,a4paper]{amsart}
\usepackage{amsmath,amssymb,amsthm,graphicx}
\usepackage{hyperref}
\usepackage{enumerate}

\numberwithin{equation}{section}

\newtheorem{thm}{Theorem}[section]
\newtheorem{cor}[thm]{Corollary}
\newtheorem{lem}[thm]{Lemma}
\newtheorem{prop}[thm]{Proposition}
{ \theoremstyle{definition}
\newtheorem{dfn}[thm]{Definition}
\newtheorem{Note}[thm]{Note}
\newtheorem{ex}[thm]{Example}
\newtheorem{rem}[thm]{Remark} }

\DeclareMathOperator{\scal}{Scal}
\DeclareMathOperator{\diam}{diam}
\DeclareMathOperator{\aire}{area}
\DeclareMathOperator{\sys}{sys}
\DeclareMathOperator{\str}{stretch}

\newcommand{\eps}{\varepsilon}
\newcommand{\ph}{\varphi}

\title{A quantitative two-systole bound for stretched positive-scalar-curvature metrics on the product of two spheres}
\author{Thomas Richard}
\date{doi:10.3842/SIGMA.2020.136; CC BY 4.0}
\begin{document}
\maketitle
\noindent\textit{Allowed standard prerequisites (editorial).}\begin{itemize}
\item Research prerequisites named in the original article. Full original bibliography retained; all source-local definitions and supporting arguments retained.
\end{itemize}\medskip
\section{Introduction}

Recall that the systole of a compact Riemannian manifold $\big(M^n,g\big)$ is the
length of the shortest non contractible loop in $\big(M^n,g\big)$. In the
middle of the 20th century, Loewner and Pu proved sharp upper bounds
on the systole of any metric on $\mathbb{T}^2$ or~$\mathbb{RP}^2$ in
term of its volume.
These were vastly generalized when in the early 80s Gromov gave similar
(non sharp) bounds for $n$-dimensional essential manifolds. (See~\cite{berger1993systoles} for the full story until 1993.)

The $k$-systole $\sys_k(g)$ of $(M^n,g)$ is the infimum of the
$k$-dimensional volume over all homologically
non trivial $k$-cycles. In general, for
$k\geq 2$, the $k$-systole of a manifold cannot be bounded by the
volume alone. In particular, Katz and Suciu showed in~\cite{katz1999volume} that one can find
metrics on $\mathbb{S}^2\times\mathbb{S}^2$ whose volume is $1$ but
whose 2-systole can be arbitrary large. (Similar examples in higher
dimension where already known to Gromov, see again~\cite{berger1993systoles}.)

One way to circumvent this is to introduce the more subtle ``stable
systoles'' (see again \cite{berger1993systoles}), another way is to
try to introduce curvature restrictions. This second route was first
considered in dimension $3$ by Bray, Brendle and Neves for manifolds
such as $\mathbb{S}^2\times\mathbb{S}^1$ with positive scalar
curvature metrics in~\cite{bray2010rigidity}. Recently, Zhu treated the case of
$\mathbb{S}^2\times\mathbb{T}^n$ ($n+2\leq 7$) in~\cite{zhu2020rigidity}, see Section~\ref{sec:results-gromov-zhu} for a
precise statement.

Here we will show how Zhu's result together with recent developments
due to Gromov gives some progress in the case of $\mathbb{S}^2\times\mathbb{S}^2$.

Let $\mathcal{S}_\ell$ be the set of embedded surfaces in
$\mathbb{S}^2\times\mathbb{S}^2$ which are in the same homology class
as $\mathbb{S}^2\times\{\ast\}$.

\begin{dfn}
 Let $g$ be a Riemannian metric on
 $\mathbb{S}^2\times\mathbb{S}^2$, the left stretch of $g$, denoted
 by $\str_{\ell}(g)$ is defined as
 \[\str_{\ell}(g)=\sup_{S_1,S_2\in\mathcal{S}_\ell}d_g\left(S_1,S_2\right).\]
\end{dfn}
\begin{ex}
 For a product metric $g=g_1\oplus g_2$ on
 $\mathbb{S}^2\times\mathbb{S}^2$, the left stretch is
 $\str_{\ell}(g)=\diam\big(\mathbb{S}^2,g_2\big)$ and is achieved by
 $S_1=\mathbb{S}^2\times\{p_1\}$ and $S_2=\mathbb{S}^2\times\{p_2\}$
 where $d_{g_2}(p_1,p_2)=\diam\big(\mathbb{S}^2,g_2\big)$.
\end{ex}
As the example above shows, the left stretch is a measure of the
1-dimensional size of the right factor of $\big(\mathbb{S}^2\times
\mathbb{S}^2,g\big)$. It is bounded above by the diameter of~$g$.
\begin{dfn}\label{dfn-str}
 Let $g$ be a Riemannian metric on
 $\mathbb{S}^2\times\mathbb{S}^2$, the left 2-systole of $g$, denoted
 by $\sys_{2,\ell}(g)$ is defined as
 \[\sys_{2,\ell}(g)=\inf_{S\in\mathcal{S}_\ell}\aire_g(S).\]
\end{dfn}
\begin{ex}
 For a product metric $g=g_1\oplus g_2$ on
 $\mathbb{S}^2\times\mathbb{S}^2$, the left 2-systole is
 $\sys_{2,\ell}(g)=\aire\big(\mathbb{S}^2,g_1\big)$ and is achieved by
 any $S=\mathbb{S}^2\times\{\ast\}$.
\end{ex}

The theorem below is the main result of the paper. It says that
positive scalar curvature metrics on $\mathbb{S}^2\times\mathbb{S}^2$
with large left stretch cannot have large left 2-systole.

\begin{thm}\label{thm-main}
 Let $g$ be a metric on $\mathbb{S}^2\times\mathbb{S}^2$ with
 $\scal_g\geq 4$.

 If $s=\str_{\ell}\big(\mathbb{S}^2\times \mathbb{S}^2,g\big) >
 \tfrac{\sqrt{3}\pi}{2}$, then $\sys_{\ell}\big(\mathbb{S}^2\times
 \mathbb{S}^2,g\big)\leq \frac{8\pi s^2}{4s^2-3\pi^2}$. Moreover there is
 an embedded $2$-sphere whose area is at most $\frac{8\pi s^2}{4s^2-3\pi^2}$.
\end{thm}

\begin{rem}
 This estimate is asymptotically sharp when $s$ goes to $+\infty$, as
 the example of the product fo two round spheres of radi
 $\frac{1}{\sqrt{k}}$ and $\frac{1}{\sqrt{2-k}}$ for $k\in[1,2)$
 shows. It remains unclear to the author wether this estimate is optimal or not. What can be shown from
 tracking down the equality case in the proofs of Theorems~\ref{thm-zhu} and~\ref{thm-gro} is that if the inequality in Theorem~\ref{thm-main} is an equality then $M$ contains an hypersurface
 $\Sigma$ such that:
 \begin{itemize}\itemsep=0pt
 \item There exists a map $\Sigma\to\mathbb{S}^2\times\mathbb{S}^1$
 with non-vanishing degre.
 \item The universal cover of $\Sigma$ is isometric to $\mathbb{S}^2\times\mathbb{R}$ with the product metric such that
 the $\mathbb{S}^2$ factor is round and has area exactly $\frac{8\pi s^2}{4s^2-3\pi^2}$.
 \end{itemize}
 Products of two round $\mathbb{S}^2$ which satisfy the properties
 above have left stretch smaller than $s$. The investigation of
 doubly warped product metrics of the form $\psi^2(r)g_{\mathbb{S}^2}+\ph^2(r){\rm d}\theta^2+ {\rm d}r^2$
 on $\mathbb{S}^2\times \mathbb{S}^2$ was inconclusive too.
\end{rem}
\begin{rem}
 Of course, if one denotes by $\mathcal{S}_r$ the set of embedded
 spheres homologous to \mbox{$\{\ast\}\times\mathbb{S}^2$}, we can define
 right stretch and systole as
 $\str_{r}(g)=\sup\limits_{S_1,S_2\in\mathcal{S}_r} d _g (S_1,S_2 )$
 and $\sys_{2,r}(g)\allowbreak =\inf\limits_{S\in\mathcal{S}_r}\aire_g(S)$ and get the
 same inequality between the right systole and the right stretch.
\end{rem}

It is currently unknown wether the $2$-systole is
bounded from above on the set
$\mathcal{M}_{\mathbb{S}^2\times\mathbb{S}^2,\scal\geq 4}$ of metrics on
$\mathbb{S}^2\times\mathbb{S}^2$ with scalar curvature greater than
$4$. Our result says that if there is no upper bound, there
must be a sequence of metrics in $\mathcal{M}_{\mathbb{S}^2\times\mathbb{S}^2,\scal\geq 4}$
whose $2$-systoles goes to infinity while their left and right
stretches stay below $\tfrac{\sqrt{3}}{2}\pi$.

The rest of the paper is organized as follows: we first review some
previously known results by Gromov and Zhu on manifolds with positive scalar curvature
which will be used to prove Theorem~\ref{thm-main}, we then prove some
topological facts that will be needed in the proof of Theorem~\ref{thm-main} and we finally give the proof of the main theorem.

\section{Positive scalar curvature inequalities by Gromov and Zhu}\label{sec:results-gromov-zhu}

We will need the following result by Jintian Zhu, which was already
alluded to in the introduction.

\begin{thm}[{\cite[Theorem 1.2]{zhu2020rigidity}}]\label{thm-zhu}
 Let $\big(M^n,g\big)$ be a closed manifold of dimension at most seven such
 that:
 \begin{itemize}\itemsep=0pt
 \item $\scal_g\geq 2$;
 \item there exists a map $F\colon M^n\to\mathbb{S}^2\times\mathbb{T}^{n-2}$
 with non-vanishing degree.
 \end{itemize}
 Then $\sys_2(g)\leq 4\pi$. More precisely one can find an
 embedded $2$-sphere $S$ such that
 $F_*([S])=\big[\mathbb{S}^2\times\{\ast\}\big]$ and whose area is less than $4\pi$.
\end{thm}
The proof uses the second variation formula of the area functional for stable minimal hypersurfaces and a
repeated symmetrization construction to reduce the problem to
$\mathbb{T}^{n-2}$-invariant metrics on
$\Sigma^2\times\mathbb{T}^{n-2}$. This method of proof was pioneered
by Fischer-Colbrie and Schoen in dimension~3 and used in
dimensions at most seven by Gromov and Lawson in~\cite{gromov1983positive}. Zhu also \mbox{studies} the equality case: under
the hypothesis of the theorem, if there is an embedded $2$-sphere~$S$
such that $F_* ([S]) = \big[\mathbb{S}^2 \times \{*\}\big]$ and whose area is~$4\pi$, then the universal cover of~$M$ is a product of the round $2$-sphere of area $4\pi$ and a euclidean space.

\begin{rem}\label{rem-zhu}
 The proof in~\cite{zhu2020rigidity} shows that if one starts with an
 $\mathbb{S}^1$-invariant metric on $M=M'\times\mathbb{S}^1$
 satisfying the hypothesis of Theorem~\ref{thm-zhu}, then one can
 find a $2$-sphere in~$M'$ whose area is at most $4\pi$.
\end{rem}

In the Spring of 2019, the author was lucky enough to attend the
lectures ``Old, New and Unknown around Scalar Curvature'' given by
Gromov at IHES. There, the author learned about the following result by
Gromov which is also proved using the Fischer-Colbrie--Schoen
symmetrization process:
\begin{thm}[{\cite[p.~2]{gromov2018metric}}]
 Let $2\leq n\leq 7$, $M=[-1,1]\times \mathbb{T}^{n-1}$ and
 $\partial_{\pm}M=\{\pm 1\}\times \mathbb{T}^{n-1}$.
 Let $g$ be a Riemannian metric on~$M$
 such that $\scal_g\geq n(n-1)$. Then
 \[d_g\big(\partial_-M,\partial_+M\big)\leq\frac{2\pi}{n}. \]
\end{thm}
The equality case is also studied. If $d_g\big(\partial_-M,\partial_+M\big)=\frac{2\pi}{n}$, then $M$ is isometric to
$\big[{-}\tfrac{\pi}{n},\tfrac{\pi}{n}\big]\times\mathbb{T}^{n-1}$
with the product metric ${\rm d}t^2 + \cos^{\tfrac{4}{n}}\big(\frac{nt}{2}\big){\rm d}x^2$
where ${\rm d}x^2$ is a flat metric on $\mathbb{T}^{n-1}$.

In his Spring 2019 lectures, Gromov gave a new proof of the previous
theorem. Instead of using the second variation of the
$(n-1)$-dimensional volume as in~\cite{gromov2018metric}, a twisted
functional is considered. Given a density $\mu\colon M^n\to\mathbb{R}$, the
$\mu$-area functional maps an open set $U\subset M$ to
$\mathcal{V}_{n-1}(\partial U)-\int_U\mu$, where $\mathcal{V}_{n-1}$
denotes the $(n-1)$-dimensional volume. Using a well chosen
density $\mu$, Gromov proved the following theorem, which is a
generalisation of the theorem above.
\begin{thm}[{\cite[Section~5.3]{gromov2019four}}]\label{thm-gro}
 Let $\big(M^n,g\big)$ $(n\leq 7)$ be a compact $n$-manifold with two boundary components $\partial_-M$
 and $\partial_+M$. Assume that
 \[\scal_g\geq \frac{4(n-1)\pi^2}{n\, d_g(\partial_-M,
 \partial_+M)^2}+\delta\]
 for some $\delta>0$. Then there exists
 \begin{itemize}\itemsep=0pt
 \item an hypersurface $\Sigma$ which separates
 $\partial_-M$ and $\partial_+M$,
 \item a positive function
 $u\colon \Sigma\to\mathbb{R}$,
 \end{itemize}
 such that the metric
 $h=g|_{\Sigma}+u^2{\rm d}t^2$ on $\Sigma\times\mathbb{R}$ has $\scal_h\geq \delta$.
\end{thm}

\section{Topological preliminaries}
\label{sec:topol-prel}

Before proving the main theorem, we establish a topological
preliminary. Let $M=\mathbb{S}^2\times\mathbb{S}^2$.

The idea is that if we remove two disjoint surfaces $S_1$ and $S_2$ in
$\mathcal{S}_\ell$ from $M$, we should be in a situation which is very
similar to removing $\mathbb{S}^2\times\{n\}$ and
$\mathbb{S}^2\times\{s\}$ from $\mathbb{S}^2\times\mathbb{S}^2$, where
$n$ and $s$ are the north and south
poles of the right factor. Thus $\tilde M=M\backslash (S_1\cup S_2)$
should look like
\[\mathbb{S}^2\times\mathbb{S}^2\backslash\big(\mathbb{S}^2\times\{n\}\cup
 \mathbb{S}^2\times\{s\}\big)=\mathbb{S}^2\times
\big(\mathbb{S}^2\backslash\{n,s\}\big)\simeq \mathbb{S}^2\times
 \mathbb{S}^1\times (-1,1).\]

However, our surfaces $S_1$ and $S_2$ from $\mathcal{S}_\ell$ may have
higher genus and thus $\tilde M$ may be more complicated.

However we are able to prove the following, which will be enough for our purpose:

\begin{prop}\label{prop-top}
 Let $S_1$ and $S_2$ be two disjoint surfaces in
 $\mathcal{S}_\ell$. Let $\Sigma\subset M$ be a connected hypersurface such
 that
 \begin{itemize}\itemsep=0pt
 \item $\Sigma$ is disjoint from $S_1$ and $S_2$,
 \item $\Sigma$ separates $S_1$ and $S_2$.
 \end{itemize}
 Then there exists a non-zero degree map $\Sigma\to\mathbb{S}^2\times\mathbb{S}^1$.
\end{prop}
This map will be the restriction of a map $F\colon\tilde M=M\backslash
(S_1\cup S_2)\to \mathbb{S}^2\times\mathbb{S}^1$.

We will first show:
\begin{lem}\label{lem-norm}
 Any $S\in \mathcal{S}_\ell$ has a trivial normal bundle.
\end{lem}
\begin{proof}
 Let $NS$ be the normal bundle to $S$ and let $e(NS)$ be the Euler
 class of~$NS$. It follows from \cite[Theorem~11.17, p.~125]{BottTu} that the Euler class of a vector bundle over a compact manifold can be computed as the intersection number between the zero section and another
 transverse section times the fundamental class of the base.

 Since $NS$ can be embedded in $M$ as a small tubular neighborhood of
 $N$, this intersection number can be computed in $M$. Since in
 $\mathcal{S}_\ell$ one can find two disjoint spheres, the
 intersection number is~0 and~$e(NS)=0$.

 Moreover the Euler class is the only obstruction for an oriented
 rank $k$ vector bundle over a~compact $k$-manifold to have a non
 vanishing section (see \cite[Proposition~3.22]{hatcher2003vector}). Hence $NS$ has a nowhere vanishing section. Since $NS$ is an
 orientable rank~2 vector bundle, it is trivial.
\end{proof}

For $\eps>0$, we denote by $S_i^\eps$ the tubular neighborhood around~$S_i$. We will choose $\eps$ small enough so that $S_1^\eps$ and
$S_2^\eps$ are disjoint regular neighborhoods.

\begin{lem}\label{lem-S2}
 The map $F_2\colon \tilde M\subset \mathbb{S}^2\times\mathbb{S}^2\to\mathbb{S}^2\times\{\ast\}$ is
 surjective in homology.
\end{lem}
\begin{proof}
 Consider a $3$-cycle $C$ in $M$ such that $\partial C= S_1-S_2$. By using
 an fine enough triangulation, we can decompose $C=C_1+\tilde C+C_2$,
 where $\tilde C$ is a $3$-cycle in $\tilde M$ and $C_i$ is a
 $3$-cycle in $S_i^\eps$. Now, $S'_1=\partial C_1-S_1$ is a $2$-cycle in
 $\tilde M$ homologous to $S_1$, hence $\pi(F_2)_*([S'_1])=\big[\mathbb{S}^2\times\{\ast\}\big]$.
\end{proof}

Set $S=S_1\cup S_2$, $S^\eps=S_1^\eps\cup S_2^\eps$ and $\tilde
S^\eps=S^\eps\backslash S$.

\begin{lem}\label{lem-H1}
 $H^1\big(\tilde M,\mathbb{Z}\big)=\mathbb{Z}$.
\end{lem}

\begin{proof}
 We will use the Mayer--Vietoris exact sequence for cohomology
 associated with the decomposition $M=S^\eps\cup \tilde M$. Note that
 $S^\eps$ is homotopy equivalent to $S$ and
 since $S_1$ and $S_2$ have trivial normal bundle, $\tilde
 S^\eps$ is homotopy equivalent to $S\times\mathbb{S}^1$.

 The Mayer--Vietoris sequence gives
 \[\cdots\to H^1(M)\to H^1\big(S^\eps\big)\oplus
 H^1\big(\tilde M\big) \to H^1\big(\tilde S^\eps\big)\to H^2(M)\to\cdots.\]
 Thus, we have
 \[\cdots\to 0\to H^1(S_1)\oplus H^1(S_2)\oplus
 H^1\big(\tilde M\big) \to H^1(S_1)\oplus\mathbb{Z}\oplus
 H^1(S_2)\oplus\mathbb{Z}\to \mathbb{Z}\oplus\mathbb{Z}\to\cdots,\]
 where the map $ H^1\big(\tilde S^\eps\big)\simeq H^1(S_1)\oplus\mathbb{Z}\oplus
 H^1(S_2)\oplus\mathbb{Z}\to H^2(M)$ is given by
 $(a,x,b,y)\mapsto (x-y,0)$. Hence $H^1\big(\tilde M\big)$ is isomorphic to~$\mathbb{Z}$.
\end{proof}

Since $H^1(\tilde M,\mathbb{Z})\simeq \mathbb{Z}$, $H^1\big(\tilde
M,\mathbb{R}\big)\simeq \mathbb{R}$ and we can find a closed $1$-form
$\alpha$ on~$\tilde M$ which is not exact. Let
$\Gamma$ be the image in $\mathbb{R}$ of the abelian group morphism:
$c\in H_1\big(\tilde M,\mathbb{Z}\big)\mapsto \int_c\alpha\in\mathbb{R}$. By a
computation similar to the proof of the previous lemma, $H_1\big(\tilde
M,\mathbb{Z}\big)=\mathbb{Z}$. Hence $\Gamma$ is a~discrete subgroup of $\mathbb{R}$. The
classical period map construction gives:
\begin{lem}\label{lem-S1}
 Let $x_0\in\tilde M$. For any curve $c$ from $x_0$ to $x$,
 $\int_c\alpha$ is well defined as an element~$\mathbb{R}/\Gamma$. This
 defines a map $F_1\colon \tilde M\to\mathbb{R}/\Gamma$ such that the induced map
 $H_1\big(\tilde M\big)\to H_1(\mathbb{R}/\Gamma) $ is an isomorphism.
\end{lem}

$F=(F_2,F_1)\colon \tilde M\to \mathbb{S}^2\times\mathbb{S}^1$ is the map we
are looking for.

\begin{proof}[Proof of Proposition~\ref{prop-top}]
First note that any connected hypersurface $\Sigma\subset M$ which separates $S_1$ and $S_2$
will be homologous to $\partial S_1^\eps$ for $\eps$ small
enough. Thus it is enough to show that the image of $[\partial
S_1^\eps]$ under $F$ is not zero in
$H_3\big(\mathbb{S}^2\times\mathbb{S}^1\big)=H_2\big(\mathbb{S}^2)\oplus
H_1(\mathbb{S}^1\big)$, which is routinely deduced from the proofs of
Lemmas~\ref{lem-S2} and~\ref{lem-H1}.
\end{proof}

We will also need an elementary fact about the connectedness of separating hypersurfaces. Before stating it let us recall that $\Sigma\subset M$ separates $K_1, K_2 \subset M$ if~$K_1$ and~$K_2$ lie in different connected components of $M\backslash\Sigma$. The result we will need is the following:

\begin{prop}\label{prop-connect}
 Let $M^n$ be a closed orientable manifold with $H_{n-1}(M) = 0$ and
 let \mbox{$\Sigma^{n-1}\!\subset\! M^n$} be a closed orientable
 hypersurface. Let $K_1$ and $K_2$ be two closed subsets of $M$ such
 that $\Sigma$ separates $K_1$ and $K_2$. Then there is a connected
 component of $\Sigma$ which separates $K_1$ and $K_2$.
\end{prop}

This is probably well known but the author has not been able to locate a proof in the litterature.

\begin{proof}Let us first remark that since $H_{n-1}(M) = 0$, any connected
 closed hypersurface $\Sigma_0\subset M$ satisfies that
 $M\backslash\Sigma_0$ has exactly two connected components whose
 boundary in $M$ is $\Sigma_0$. The fact that $M\backslash\Sigma_0$ has at most two
 connected components follows from the fact $\Sigma_0$ is two sided
 (since both $\Sigma_0$ and $M$ are orientable). If
 $M\backslash\Sigma_0$ was connected, there would be a closed loop
 $\gamma$ which intersects $\Sigma_0$ exactly once, hence the
 intersection number of $[\gamma]\in H_1(M)$ and $[\Sigma_0]\in
 H_{n-1}(M)$ would be non zero, which is absurd since $H_{n-1}(M)=0$.

 We will now prove the following result: let $\tilde\Sigma$ and $\Sigma_0$ be two disjoint hypersurfaces such
 that:
\begin{itemize}\itemsep=0pt
\item $\Sigma_0$ is connected.
\item $\tilde\Sigma\cup\Sigma_0$ separates $K_1$ and $K_2$.
\item $\Sigma_0$ does not separate $K_1$ and $K_2$.
\end{itemize}
Then $\tilde\Sigma$ separates $K_1$ and $K_2$.

Once we have proved this, the proposition follows since we can remove
all the non-separating components form $\Sigma$ until every connected
component of $\Sigma$ separates $K_1$ and $K_2$.

In order to prove the result above, let $U_0$ be the connected
component of $M\backslash\Sigma_0$ which contains~$K_1$. Since
$\Sigma_0$ does not separate $K_1$ and~$K_2$, $U_0$ also contains
$K_2$. Now let $\tilde U$ be the connected component of
$M\backslash\tilde\Sigma$ which contains~$K_1$.

Since the boundaries of $U_0$ and $\tilde U$ are disjoint and both
$U_0$ and $\tilde U$ contain $K_1$ either
$U_0 \subset \tilde U$ or $\tilde U \subset U_0$. In the first case,
$U_0$ will be a connected component of $M\backslash\Sigma$ which
contains~$K_1$ and~$K_2$, which is impossible since $\Sigma$ separates~$K_1$ and~$K_2$.

Hence $\tilde U \subset U_0$, which implies that $\tilde U$ is one of
the connected components of $M\backslash\Sigma$. Since~$\Sigma$ separates
$K_1$ and $K_2$, $K_2$ cannot be included in~$\tilde U$. Hence~$\tilde\Sigma$ separates
$K_1$ and~$K_2$.
\end{proof}
\section{Proof of the main theorem}
\label{sec:proof}

We are now ready to prove Theorem \ref{thm-main}.

\begin{proof}[Proof of Theorem \ref{thm-main}]
 Our assumption is that $s=\str_{\ell}\big(\mathbb{S}^2\times \mathbb{S}^2,g\big) >
 \tfrac{\sqrt{3}\pi}{2}$. Thus, for any $\eps>0$ we can find two surfaces
 $S_1$ and $S_2$ in $\mathcal{S}_\ell$ such that $d_g(S_1,S_2)=
 s-\eps$. Assume that $\eps>0$ is such that
 \begin{itemize}\itemsep=0pt
 \item $s-3\eps>\frac{\sqrt{3}\pi}{2}$;
 \item the tubular neighborhoods $S_1^\eps$ and $S_2^\eps$ of radius
 $\eps$ around $S_1$ and $S_2$ are disjoint and regular.
 \end{itemize}
 Then if we set $\tilde M=M\backslash \big(S_1^\eps\cup S_2^\eps\big)$, its
 two boundary components $\partial_-\tilde M=\partial S_1^\eps$ and
 $\partial_+\tilde M=\partial S_2^\eps$ satisfy $d_g\big(\partial_-\tilde
 M, \partial_+\tilde M\big)= s-3\eps$.

 Set $\delta=4-\frac{3\pi^2}{(s-3\eps)^2}$ and note that $\delta>0$. Then we have
 \[\scal_g\geq 4=\frac{3\pi^2}{d_g\big(\partial_-\tilde
 M, \partial_+\tilde M\big)^2}+\delta.\]

 This is exactly what we need to apply Gromov's Theorem~\ref{thm-gro}. Hence we get an hypersurface $\Sigma\subset\tilde M$
 which separates $S_1$ and $S_2$ such that
 $\Sigma\times\mathbb{S}^1$ admits a metric of the form
 $h=g_\Sigma+u^2{\rm d}t^2$ for some $u\colon M\to\mathbb{R}$ such that
 $\scal_h\geq\delta$. The hypersurface coming the minimization
 process in the proof of Theorem~\ref{thm-gro} may be disconnected,
 however by Proposition~\ref{prop-connect} we can replace~$\Sigma$ by
 one of its connected components which separates~$S_1$ and~$S_2$.

 Moreover, by Proposition~\ref{prop-top}, there exists a non-zero
 degree map $\Sigma\to\mathbb{S}^2\times\mathbb{S}^1$. Hence there is
 a non-zero degree map
 $\Sigma\times\mathbb{S}^1\to\mathbb{S}^2\times\mathbb{T}^2$ and we
 can apply Zhu's Theorem~\ref{thm-zhu} and Remark~\ref{rem-zhu} to
 show that one can find an embedded $2$-sphere $S$ in $\Sigma$ whose
 area is at most $\frac{8\pi}{\delta}=\frac{8\pi
 (s-3\eps)^2}{4(s-3\eps)^2-3\pi^2}$.

 Since $\Sigma$ is isometrically embedded in
 $M$, $S$ embeds in $M$ with area at most $\frac{8\pi
 (s-3\eps)^2}{4(s-3\eps)^2-3\pi^2}$. Moreover, by the properties of the map $F$
 built in Section~\ref{sec:topol-prel}, the surface $S$ of $M$
 belongs to $\mathcal{S}_\ell$. Since $\eps>0$ can be as small as one
 wants we get the results.
\end{proof}


\subsubsection*{Acknowledgements}

The author would like to thank G.~Besson, S.~Maillot and S.~Sabourau for helpful discussions. He also thanks Jintian Zhu and the anonymous referees for pointing inaccuracies in a previous version of this work. The author is supported by the grant ANR-17-CE40-0034 of the French National Research Agency ANR (Project CCEM).

\pdfbookmark[1]{References}{ref}

\begin{thebibliography}{99}
\footnotesize\itemsep=0pt

\bibitem{berger1993systoles}
Berger M., Systoles et applications selon {G}romov, \textit{Ast\'erisque}
 \textbf{216} (1993), 771, 279--310.

\bibitem{BottTu}
Bott R., Tu L.W., Differential forms in algebraic topology, \textit{Graduate
 Texts in Mathematics}, Vol.~82, \href{https://doi.org/10.1007/978-1-4757-3951-0}{Springer-Verlag}, New York~-- Berlin, 1982.

\bibitem{bray2010rigidity}
Bray H., Brendle S., Neves A., Rigidity of area-minimizing two-spheres in
 three-manifolds, \href{https://doi.org/10.4310/CAG.2010.v18.n4.a6}{\textit{Comm. Anal. Geom.}} \textbf{18} (2010), 821--830,
 \href{https://arxiv.org/abs/1002.2814}{arXiv:1002.2814}.

\bibitem{gromov2018metric}
Gromov M., Metric inequalities with scalar curvature, \href{https://doi.org/10.1007/s00039-018-0453-z}{\textit{Geom. Funct.
 Anal.}} \textbf{28} (2018), 645--726, \href{https://arxiv.org/abs/1710.04655}{arXiv:1710.04655}.

\bibitem{gromov2019four}
Gromov M., Four lectures on scalar curvature, \href{https://arxiv.org/abs/1908.10612}{arXiv:1908.10612}.

\bibitem{gromov1983positive}
Gromov M., Lawson Jr. H.B., Positive scalar curvature and the {D}irac operator
 on complete {R}iemannian manifolds, \href{https://doi.org/10.1007/BF02953774}{\textit{Inst. Hautes \'Etudes Sci. Publ.
 Math.}} \textbf{58} (1983), 83--196.

\bibitem{hatcher2003vector}
Hatcher A., Vector bundles and K-theory, available at
 \url{http://pi.math.cornell.edu/~hatcher/}.

\bibitem{katz1999volume}
Katz M.G., Suciu A.I., Volume of {R}iemannian manifolds, geometric
 inequalities, and homotopy theory, in Tel {A}viv {T}opology {C}onference:
 {R}othenberg {F}estschrift (1998), \textit{Contemp. Math.}, Vol.~231, \href{https://doi.org/10.1090/conm/231/03357}{Amer.
 Math. Soc.}, Providence, RI, 1999, 113--136, \href{https://arxiv.org/abs/math.DG/9810172}{arXiv:math.DG/9810172}.

\bibitem{zhu2020rigidity}
Zhu J., Rigidity of area-minimizing {$2$}-spheres in {$n$}-manifolds with
 positive scalar curvature, \href{https://doi.org/10.1090/proc/15033}{\textit{Proc. Amer. Math. Soc.}} \textbf{148}
 (2020), 3479--3489, \href{https://arxiv.org/abs/1903.05785}{arXiv:1903.05785}.

\end{thebibliography}
\end{document}
